% Generated by report/make_restated.py from theory/B3_prop59_tightness.tex: do not edit (edit the note and rerun).
\begin{restated}{Remark}{rem:B3-R2}{B3}{0}{the constant $R^2$ for logistic regression}
For $f_i(\theta)=\ell(y_ix_i^\top\theta)+\frac\lambda2\norm\theta^2$ with $\ell(u)=\log(1+e^{-u})$, $y_i\in\{-1,1\}$, one has $\ell''(u)=\sigma(u)\sigma(-u)\in(0,\frac14]$ where $\sigma(u)=1/(1+e^{-u})$, so $f_i''(\theta)=\ell''(y_ix_i^\top\theta)\,x_ix_i^\top+\lambda I\preceq\big(\frac{\norm{x_i}^2}{4}+\lambda\big)I$. Hence $f_i$ is $L_i$-smooth with $L_i=\norm{x_i}^2/4+\lambda$, and the smallest constant valid for all $i$ is
\[
R^2:=\max_{i}\Big(\frac{\norm{x_i}^2}{4}+\lambda\Big).
\]
With features bounded by $R_x$ in $\ell_2$-norm, the book's constant $R^2:=R_x^2$ (``features bounded by $R$ in $\ell_2$-norm'') is therefore a valid smoothness constant whenever $\lambda\le\frac34R_x^2$, but a loose one: it ignores the factor $\frac14$ of the logistic second derivative (and $\lambda$), so that it exceeds $\max_iL_i$ by a factor close to $4$ when $\lambda\ll R_x^2$. Since every constant of the proof is expressed through $R^2$ ($\gamma=1/(4R^2)$, $\Delta=3/(16R^4)$, the rate $3\mu/(16R^2)$), using the loose value would divide the step and the guaranteed rate by about $4$; the experiments must use $R^2=\max_i(\norm{x_i}^2/4+\lambda)$. The global strong convexity constant is $\mu=\lambda$ (the logistic part is convex), which is much smaller than the curvature actually met along the trajectory when $\lambda$ is small compared with the curvature of the data term: in the experiments of Section~\ref{sec:B3-measure}, the ratio $\muEff/\mu$ of Definition~\ref{def:B3-ratios} (geometric means over passes 2--20) lies between $1.1$ and $2.3$ for $\kappa\le10^2$ and reaches $81$ at $\kappa=10^5$; see Section~\ref{sec:B3-traj}.
\end{restated}
